\subsection{三高 core 間の単位枝：小分母を持つ二種類の境界}
\label{nu24:unit-three-high}

標準五次の単位枝 $ij$ の両端が正の三高 core である場合，低行の正成分は
異なる第三 core を指す．共通に負となる唯一の位置を $c$ とし，低行の絶対値を
$s,s'$ とすると，$n_{ij;c}=1+s+s'\le5$ である．
第三 core の pair を課すと $s=s'=1$ は不可能になり，node 配分は
\[
 \{5,0,0\}\quad\text{または}\quad\{4,1,0\}
\]
だけになる．これは分母の探索上限を使わない必要条件である．

\begin{center}\small
\begin{tabular}{c|c|c}
共通負位置の node&隣接する二高分母&その第三 core の pair\\\hline
5&$(3,7),(5,5),(7,3)$&負\\
4&$(3,3)$&負\\
4&$(3,2R+1),(2R+1,3)$，$R\ge1$&正
\end{tabular}
\end{center}
従って各単位枝は，隣接するどこかに三または五の小分母を強制する．
この結論は $G_{5,14}$ の両単位枝だけでなく，$G_{5,16}$ の唯一の単位枝にも使える．

五nodeの場合，共通負位置の両低 column の絶対値の和は四である．
他の二位置の node は零で，正 column と反対側の負 column の絶対値の差は
それぞれ零または一である．対向低 $D$ 相殺がその和を二にするため，
両方とも一となる．従って他の二つの第三 core の pair は正で，対応する
高辺の node は零である．
$G_{5,14}$ では五node位置はもう一方の単位枝の端点に限られ，四高 core にはならない．

さらに対向低 $D$ 相殺を使うと，任意の単位枝の三つの外側 pair の少なくとも一つは正である．
全て負と仮定した場合は，完全三角形表から一方の低データが
\[
 \text{row}=(-2,-1,2),\qquad\text{column}=(-1,-2,2r),\qquad r\ge1
\]
に強制される．局所完成の残余は
$(3A+6r+1)(A+6r+1)=0$ であり，非整数または零 entry を強制して不可能になる．

\paragraph{再利用する完全な三角形表．}
反対符号位置では，低 row を $u,-j_0$，column を $U,-T$ と書くと
$j_0=u+n-1$，$n\in\{0,1,2\}$ である．
第三 core の pair は
\[
 ((u\alpha-1)U+\alpha)((j_0\beta+1)T+\beta)
   =(uU+1)(j_0T+1)-1
\]
となる．四通常族は
\[
\begin{array}{c|c|c}
 n&(\alpha,\beta)&尺度条件\\\hline
 0&(1,1)&T=U-1\\
 1&(1,1)&U=1\\
 1&(-1,-1)&T=1\\
 2&(-1,-1)&T=U+1
\end{array}
\]
である．正の分母は奇数なので，最初と最後の族では $u,U$ の偶奇が逆，
中二族では $u$ が偶数である．
小分母三・五を含む四種類の例外もあり，公開稿の完全表から必ず併用する．
通常族だけで尺度を打ち切ってはいけない．
両端が正の三高である実際の単位枝では，共通負位置の分類によって $n=2$ 自体が除かれるが，
単一三角形の補題としては第四通常族を残しておく．

\paragraph{次の入口．}
$G_{5,14}$ の両単位枝は $\{4,1,0\}$ に限られ，中央の四高 core の正 matching も固定された．
次は二つの node 四の共通負位置，隣接分母三，固定中央 matching を同時に合わせる．
五node sector を引き続き調べるのは $G_{5,16}$ であり，二単位枝を使う排除を移さない．
四nodeの場合には自由尺度 $R$ が残るので，有限表を分母全体の有限性と混同しない．
両端の完全行和・pair products，対向高 $D$，全高辺の node 総数を復元時に戻す．

公開稿は \nolinkurl{unit_three_high_triangle_reduction_v24.tex}，
検算は \nolinkurl{verify_unit_three_high_triangles_v24.py}．
共通負位置の九十六負符号場合は整数因数分解で完全に検証され，二十四正符号場合も記録する．
反対符号族の無上界分類は本文の符号・整数間隙の証明に依存し，
有界回帰は実装上の取り違えを検出する補助としてのみ扱う．
